A metodologia adotada neste trabalho está organizada em três etapas, que permitem analisar o impacto da substituição dos Online Codes pelo Raptor R10 no HySAIL quanto à integridade e à disponibilidade dos dados e à eficiência dos processos de codificação e decodificação. Na primeira etapa, realizou-se uma revisão da literatura que fundamenta a proposta. Na segunda, avalia-se se as provas de integridade e disponibilidade do HySAIL permanecem válidas após a substituição. Na terceira, compara-se experimentalmente o desempenho do R10 com o dos Online Codes e com o do RaptorQ.

\section{Revisão da Literatura}
A revisão da literatura abrangeu os Online Codes, os Raptor Codes, nas variantes R10 e RaptorQ, e as soluções que empregam códigos de apagamento para garantir a integridade e a disponibilidade de dados armazenados. Seu objetivo foi compreender o estado da arte e identificar as lacunas que justificam a substituição proposta.

A busca foi conduzida de forma sistemática nas bases Scopus, Web of Science, Google Scholar e IEEE Xplore, com as palavras-chave ``Online Codes'', ``Raptor Codes'', ``RaptorQ Codes'', ``proof of retrievability'', ``erasure codes'', ``'' e ``digital storage'', além de seus sinônimos e variações em português e inglês. Para assegurar a atualidade dos resultados, foram considerados artigos publicados nos últimos oito anos. A análise dos trabalhos selecionados permitiu identificar as lacunas e oportunidades que fundamentam a proposta apresentada neste trabalho.

\section{Provas de Integridade e Disponibilidade}
O HySAIL fundamenta-se em provas de integridade e disponibilidade, que garantem que os dados armazenados permanecem corretos e acessíveis mesmo diante de falhas ou de ataques maliciosos. Como essas provas foram formuladas considerando o uso de Online Codes, é necessário verificar se continuam válidas quando o Raptor R10 é adotado. Essa verificação compreenderá uma análise teórica das propriedades do R10 relacionadas à integridade e à disponibilidade, complementada por experimentos que simulem falhas e ataques, a fim de confirmar que as garantias originais são preservadas.

\section{Avaliação do Raptor R10}
Além da preservação das garantias de segurança, é necessário avaliar o impacto da substituição sobre a eficiência de codificação e decodificação e sobre a robustez frente a perdas de símbolos. Para isso, serão realizados experimentos que meçam, para o R10, os Online Codes e o RaptorQ, o \emph{overhead} de recepção, o tempo de codificação e de decodificação e a taxa de recuperação dos símbolos de origem, permitindo uma análise comparativa entre as duas abordagens.
